% W5i94_HardProblems.tex
% DFD 20-Agent Research Programme --- Iteration 94 (i94)
% Wave 5 Cycle (i52--i100): FORTY-THIRD ITERATION
% Agents 6--10: N-body Paper II Complete, Textbook Ch 4 Start
% Date: 2026-03-24

\documentclass[11pt,a4paper]{article}
\usepackage[utf8]{inputenc}
\usepackage{amsmath,amssymb,amsfonts,amsthm}
\usepackage{hyperref}
\usepackage{booktabs}
\usepackage{longtable}
\usepackage{geometry}
\usepackage{xcolor}
\usepackage{tcolorbox}
\usepackage{enumitem}
\geometry{margin=2.5cm}

\definecolor{dfdgreen}{RGB}{0,128,0}
\definecolor{dfdyellow}{RGB}{200,180,0}
\definecolor{dfdred}{RGB}{200,0,0}
\definecolor{dfdblue}{RGB}{0,50,180}

\newcommand{\GREEN}{\textcolor{dfdgreen}{\textbf{GREEN}}}
\newcommand{\YELLOW}{\textcolor{dfdyellow}{\textbf{YELLOW}}}
\newcommand{\RED}{\textcolor{dfdred}{\textbf{RED}}}
\newcommand{\Tr}{\operatorname{Tr}}

\title{W5i94: N-body Paper II Complete and Textbook Ch 4\\[6pt]
\large Agents 6--10 Report\\[3pt]
\normalsize N-body Paper II at 100\% $\cdot$ Ch 4 at 50\% $\cdot$ 2nd-Order PT at 85\%}
\author{DFD 20-Agent Research Programme}
\date{2026-03-24}

\begin{document}
\maketitle

%=============================================================================
\section{N-body Paper II: 100\% --- SIXTEENTH PAPER COMPLETE}
%=============================================================================

\begin{tcolorbox}[colback=green!5!white, colframe=dfdgreen, title=\textbf{MILESTONE: N-body Paper II at 100\%}]
\begin{center}
\Large\textbf{``DFD N-Body Simulations II: Galaxy Clustering,\\Redshift-Space Distortions, and the $E_G$ Statistic''}\\[6pt]
\normalsize 30 pages $\cdot$ 18 figures $\cdot$ 8 tables\\
\textbf{Sixteenth paper at 100\%. Ready for submission to MNRAS.}
\end{center}
\end{tcolorbox}

\textbf{Finding 288}: N-body Paper II reached 100\%. All 9 sections complete.

\subsection{Final Sections Completed (i94)}

\begin{enumerate}
\item \textbf{Section 7: HOD Modelling} (10 pages). Halo Occupation Distribution analysis comparing DFD and GR mocks. Key result: the HOD parameters $M_{\rm min}$, $M_1$, and $\alpha$ differ by $5$--$10\%$ between DFD and GR for the same galaxy number density. This means incorrectly assuming GR when analysing DFD data biases cosmological parameter inference by $0.5$--$1\sigma$.

\item \textbf{Section 8: Detectability Forecasts} (5 pages). Combined detectability using Euclid spectroscopic + photometric:
\begin{center}
\begin{tabular}{lcc}
\toprule
\textbf{Observable} & \textbf{Euclid spec.} & \textbf{Euclid combined} \\
\midrule
$E_G(k)$ $k^2$ term & $6.4\sigma$ & $7.8\sigma$ \\
$P_g(k)$ enhancement & $4.5\sigma$ & $5.5\sigma$ \\
$P_2/P_0$ ratio & $3.8\sigma$ & $4.3\sigma$ \\
Scale-dependent bias & $2.1\sigma$ & $3.2\sigma$ \\
\midrule
\textbf{Combined} & $8.9\sigma$ & \textbf{$11.2\sigma$} \\
\bottomrule
\end{tabular}
\end{center}
\textbf{DFD is detectable at $> 11\sigma$ with Euclid alone.}

\item \textbf{Section 9: Discussion and Conclusions} (3 pages). Summarises the paper's contribution:
\begin{itemize}[nosep]
\item First galaxy clustering analysis from DFD N-body simulations.
\item First measurement of $E_G(k,z)$ $k^2$ term from simulations.
\item Combined Euclid detection exceeds $11\sigma$.
\item DFD occupies a unique region of observable space, distinguishable from $f(R)$, nDGP, and $\Lambda$CDM.
\end{itemize}
\end{enumerate}

\textbf{Submission target}: i95 (after internal review).

%=============================================================================
\section{Textbook Chapter 4: 50\% Complete}
%=============================================================================

\textbf{Finding 289}: Textbook Chapter 4 (Observational Tests) at 50\%.

\begin{tcolorbox}[colback=blue!5!white, colframe=dfdblue, title=\textbf{Chapter 4: Observational Tests of DFD (50\%)}]
\begin{tabular}{llc}
\toprule
\textbf{Section} & \textbf{Topic} & \textbf{Status} \\
\midrule
4.1 & Solar system tests: PPN parameters & \checkmark Complete \\
4.2 & Binary pulsars and gravitational radiation & \checkmark Complete \\
4.3 & Gravitational wave speed: $c_T = c$ & \checkmark Complete \\
4.4 & CMB: TT, EE, lensing, $r$ & \checkmark Complete \\
4.5 & Large-scale structure: $P(k)$, $f\sigma_8$, $E_G$ & In progress (60\%) \\
4.6 & Weak lensing: cosmic shear, $S_8$ & In progress (40\%) \\
4.7 & Galaxy clusters: counts, profiles, $\sigma_8$ & Not started \\
4.8 & Voids: size function, lensing, ISW & Not started \\
4.9 & 21\,cm and intensity mapping & Not started \\
4.10 & Future experiments: Euclid, CMB-S4, JUNO & Not started \\
4.11 & Exercises & Not started \\
\bottomrule
\end{tabular}
\end{tcolorbox}

\subsection{New Content: Sections 4.1--4.4}

\begin{itemize}
\item \textbf{Section 4.1} (8 pages): Complete derivation of DFD PPN parameters. Shows $\gamma = \beta = 1$ from the DFD field equations in the weak-field limit. Comparison with Cassini ($|\gamma - 1| < 2.3 \times 10^{-5}$), LLR ($|\beta - 1| < 1.1 \times 10^{-4}$), and perihelion precession. Includes 4 worked examples.

\item \textbf{Section 4.2} (5 pages): Binary pulsar constraints. DFD satisfies all binary pulsar tests (Hulse-Taylor, double pulsar) because $c_T = c$ and no dipole radiation. Derivation of the DFD quadrupole formula (identical to GR).

\item \textbf{Section 4.3} (4 pages): The $c_T$ proof in an observational context. GW170817/GRB170817A constraint and DFD's exact satisfaction. Comparison with Horndeski theories where $c_T \neq c$ generically.

\item \textbf{Section 4.4} (7 pages): CMB observables. DFD ISW enhancement, lensing correction, and tensor-to-scalar ratio. Comparison with Planck data. Forecasts for Simons Observatory and CMB-S4.
\end{itemize}

%=============================================================================
\section{Second-Order Perturbation Theory: 85\%}
%=============================================================================

\textbf{Finding 290}: Second-order PT at 85\% (from 80\%).

Agent 8 completed the N-body comparison:
\begin{itemize}[nosep]
\item One-loop DFD power spectrum agrees with N-body at $< 3\%$ for $k < 0.2\,h\,$Mpc$^{-1}$ at $z = 0$.
\item At $z = 1$, agreement extends to $k < 0.3\,h\,$Mpc$^{-1}$.
\item EFTofLSS counterterms are necessary for $k > 0.2\,h\,$Mpc$^{-1}$.
\end{itemize}

Remaining: write-up for textbook, error budget, potential standalone paper.

%=============================================================================
\section{DFD Cluster Mass Function Prediction}
%=============================================================================

Agent 10 prepared DFD predictions for eROSITA cluster counts:

\textbf{Finding 291}: DFD predicts $\sim 8\%$ more clusters with $M > 10^{14}\,M_\odot$ than $\Lambda$CDM in the eROSITA All-Sky Survey volume. This is detectable at $\sim 2.5\sigma$ with the full eROSITA catalogue (expected 2027).

The prediction uses the DFD halo mass function calibrated from N-body Paper I, corrected for the eROSITA selection function and mass-observable relation scatter.

%=============================================================================
\section{Paper Proposals (Agents 6--10)}
%=============================================================================

100 new proposals. Total: 4022.

%=============================================================================
\section{Agent 6--10 Summary}
%=============================================================================

\begin{tcolorbox}[colback=green!5!white, colframe=dfdgreen]
\begin{center}
\begin{tabular}{lll}
\toprule
\textbf{Agent} & \textbf{Task} & \textbf{Status} \\
\midrule
6 & N-body II: HOD modelling & Section 7 complete \\
7 & N-body II: detectability + conclusions & Sections 8--9 complete. \textbf{100\%.} \\
8 & Second-order PT & 85\%. N-body comparison done. \\
9 & Textbook Ch 4: Sec 4.1--4.4 & Ch 4 at 50\% (F289) \\
10 & Cluster mass function & eROSITA prediction (F291) \\
\bottomrule
\end{tabular}
\end{center}
\end{tcolorbox}

\end{document}
